\section{Introduction}

Ballot-polling risk-limiting audits are usually presented as a statistical test:
sample ballots at random, update a likelihood ratio, and stop when the reported
winner is supported at the chosen risk limit. As with sampler replay (V.7), a
large part of public trust sits one layer below that test. An observer should be
able to ask whether the declared per-draw statistic and stop decision actually
follow from the reported totals, the risk limit, and the sequence of observed
ballots.

RCOUNT therefore treats BRAVO as a replay contract rather than a number to be
trusted. The verifier does not re-run a vendor audit tool or accept an exported
risk value. It reconstructs each observation from the public audit record and
recomputes the BRAVO likelihood ratio in integer arithmetic, so the published
transcript can be regenerated from public evidence.

This paper covers the BRAVO ballot-polling slice. It does not claim that any
particular election was audited correctly, nor that BRAVO is the right method
for a given jurisdiction. It claims that, when an artifact contains ballot-polling
observations and reported two-candidate totals, the resulting stop-or-continue
transcript is reproducible.
